% !TEX root = ../../ctfp-print.tex

\lettrine[lhang=0.17]{如}{果我们将} 自函子解释为定义表达式的方式，则代数允许我们对其进行求值，而单子允许我们构造和操作这些表达式。通过将代数与单子结合，我们不仅获得了大量功能，还能回答一些有趣的问题。

其中一个问题是单子与伴随关系之间的关联。正如我们所见，每个伴随关系 \hyperref[monads-categorically]{定义一个单子}（和一个余单子）。问题是：每个单子（余单子）是否都能由某个伴随关系导出？答案是肯定的。存在一整族伴随关系可以生成给定的单子。我将向你展示其中两个伴随关系。

\begin{figure}[H]
  \centering
  \includegraphics[width=0.25\textwidth]{images/pigalg.png}
\end{figure}

\noindent
让我们先回顾定义。单子是一个配备两个满足相容条件的自然变换的自函子 $m$。这些变换在对象 $a$ 处的分量为：
\begin{align*}
  \eta_a & \Colon a \to m\ a         \\
  \mu_a  & \Colon m\ (m\ a) \to m\ a
\end{align*}
该自函子的一个代数则是选择特定对象——载体 $a$——以及态射：
\[\mathit{alg} \Colon m\ a \to a\]
首先需要注意的是，代数的方向与 $\eta_a$ 相反。直观地说，$\eta_a$ 从类型为 $a$ 的值创建了一个平凡表达式。第一个相容条件要求代数与单子兼容，即用载体为 $a$ 的代数对该表达式求值能返回原始值：
\[\mathit{alg} \circ \eta_a = \id_a\]
第二个条件源于对双重嵌套表达式 $m\ (m\ a)$ 存在两种求值方式：我们可以先应用 $\mu_a$ 展平表达式再使用代数求值器；或者先应用提升后的求值器计算内部表达式，再对结果应用求值器。我们期望这两种策略等价：
\[\mathit{alg} \circ \mu_a = \mathit{alg} \circ m\ \mathit{alg}\]
此处 \code{m alg} 是通过函子 $m$ 提升 $\mathit{alg}$ 得到的态射。以下交换图描述了这两个条件（为后续推导方便将 $m$ 替换为 $T$）：

\begin{figure}[H]
  \centering
  \begin{subfigure}
    \centering
    \begin{tikzcd}[column sep=large, row sep=large]
      a \arrow[rd, equal] \arrow[r, "\eta_a"]
      & Ta \arrow[d, "\mathit{alg}"] \\
      & a
    \end{tikzcd}
  \end{subfigure}
  \hspace{1cm}
  \begin{subfigure}
    \centering
    \begin{tikzcd}[column sep=large, row sep=large]
      T(Ta) \arrow[r, "T\ \mathit{alg}"] \arrow[d, "\mu_a"]
      & Ta \arrow[d, "\mathit{alg}"] \\
      Ta \arrow[r, "\mathit{alg}"]
      & a
    \end{tikzcd}
  \end{subfigure}
\end{figure}

\noindent
我们也可以用 Haskell 表达这些条件：

\src{snippet01}
来看一个小例子。列表自函子的一个代数包含某个类型 \code{a} 和一个从 \code{a} 列表生成 \code{a} 的函数。我们可以用 \code{foldr} 表达这个函数，令元素类型和累加器类型都等于 \code{a}：

\src{snippet02}
这个特定代数由一个双参数函数（记作 \code{f}）和一个值 \code{z} 确定。列表函子恰好也是一个单子，其 \code{return} 将值转换为单元素列表。代数（此处为 \code{foldr f z}）与 \code{return} 的复合将 \code{x} 映射为：

\src{snippet03}
其中 \code{f} 的作用以中缀形式书写。若以下相容条件对任意 \code{x} 成立，则该代数与单子兼容：

\src{snippet04}
若将 \code{f} 视为二元运算符，此条件表明 \code{z} 是右单位元。

第二个相容条件作用于列表的列表。\code{join} 的作用是连接各子列表，我们可以折叠合并后的列表；另一方面，也可以先分别折叠各子列表，再折叠合并后的列表。若将 \code{f} 视为二元运算符，此条件表明该运算是可结合的。当 \code{(a, f, z)} 构成幺半群时，这些条件显然成立。

\section{T-代数}

由于数学家习惯将单子记为 $T$，他们称与其兼容的代数为 T-代数。给定范畴 $\cat{C}$ 中单子 $T$ 的 T-代数构成一个称为 Eilenberg-Moore 范畴的范畴，常记作 $\cat{C}^T$。该范畴中的态射是代数同态，即与之前 F-代数定义相同的同态。

T-代数由载体对象和求值器组成的对 $(a, f)$ 构成。存在一个明显的遗忘函子 $U^T$ 从 $\cat{C}^T$ 到 $\cat{C}$，它将 $(a, f)$ 映射为 $a$。同时它也将 T-代数同态映射为范畴 $\cat{C}$ 中载体对象间的相应态射。你或许记得在讨论伴随关系时，遗忘函子的左伴随被称为自由函子。

$U^T$ 的左伴随记为 $F^T$。它将范畴 $\cat{C}$ 中的对象 $a$ 映射为 $\cat{C}^T$ 中的自由代数。该自由代数的载体是 $T a$，其求值器是从 $T (T a)$ 到 $T a$ 的态射。由于 $T$ 是单子，我们可以用单子的 $\mu_a$（Haskell 中的 \code{join}）作为求值器。

我们仍需验证这是 T-代数。为此需满足两个相容条件：
\begin{align*}
  \mathit{alg} & \circ \eta_{Ta} = \id_{Ta}     \\
  \mathit{alg} & \circ \mu_a = \mathit{alg} \circ T\ \mathit{alg}
\end{align*}
但若将代数替换为 $\mu$，这些正是单子定律。

如你回忆，每个伴随关系定义一个单子。事实证明，$F^T$ 与 $U^T$ 之间的伴随关系定义的正是用于构建 Eilenberg-Moore 范畴的那个单子 $T$。由于这种构造适用于每个单子，我们得出结论：每个单子都可以由某个伴随关系生成。稍后我将展示另一个生成相同单子的伴随关系。

计划如下：首先我将证明 $F^T$ 确实是 $U^T$ 的左伴随，通过定义该伴随关系的单位和余单位并验证相应的三角恒等式。然后我将证明该伴随关系生成的单子确实是我们原来的单子。

该伴随关系的单位是自然变换：
\[\eta \Colon I \to U^T \circ F^T\]
计算该变换在 $a$ 处的分量：恒等函子给出 $a$，自由函子产生自由代数 $(T a, \mu_a)$，遗忘函子将其简化为 $T a$。整体得到从 $a$ 到 $T a$ 的映射。我们将简单地使用单子 $T$ 的单位作为该伴随关系的单位。

观察余单位：
\[\varepsilon \Colon F^T \circ U^T \to I\]
计算它在某个 T-代数 $(a, f)$ 处的分量。遗忘函子忘记 $f$，自由函子产生对 $(T a, \mu_a)$。因此为了定义余单位 $\varepsilon$ 在 $(a, f)$ 处的分量，我们需要 Eilenberg-Moore 范畴中的适当态射，即 T-代数同态：
\[(T a, \mu_a) \to (a, f)\]
这样的同态应将载体 $T a$ 映射到 $a$。让我们重新启用被遗忘的求值器 $f$。这次我们将它用作 T-代数同态。事实上，使 $f$ 成为 T-代数的交换图也可重新解释为显示它是 T-代数同态：

\begin{figure}[H]
  \centering
  \begin{tikzcd}[column sep=large, row sep=large]
    T(Ta) \arrow[r, "T f"] \arrow[d, "\mu_a"]
    & Ta \arrow[d, "f"] \\
    Ta \arrow[r, "f"]
    & a
  \end{tikzcd}
\end{figure}

\noindent
因此我们定义余单位自然变换 $\varepsilon$ 在 $(a, f)$（T-代数范畴中的对象）处的分量为 $f$。

要完成伴随关系还需验证单位和余单位满足三角恒等式：

\begin{figure}[H]
  \centering
  \begin{subfigure}
    \centering
    \begin{tikzcd}[column sep=large, row sep=large]
      Ta \arrow[rd, equal] \arrow[r, "T \eta_a"]
      & T(Ta) \arrow[d, "\mu_a"] \\
      & Ta
    \end{tikzcd}
  \end{subfigure}%
  \hspace{1cm}
  \begin{subfigure}
    \centering
    \begin{tikzcd}[column sep=large, row sep=large]
      a \arrow[rd, equal] \arrow[r, "\eta_a"]
      & Ta \arrow[d, "f"] \\
      & a
    \end{tikzcd}
  \end{subfigure}
\end{figure}

\noindent
第一个恒等式因单子 $T$ 的单位律成立。第二个是 T-代数 $(a, f)$ 的基本性质。

我们已建立两个函子构成伴随关系：
\[F^T \dashv U^T\]
每个伴随关系产生一个单子。往返函子
\[U^T \circ F^T\]
是范畴 $\cat{C}$ 中生成相应单子的自函子。观察它在对象 $a$ 上的作用：$F^T$ 创建的自由代数是 $(T a, \mu_a)$，遗忘函子 $U^T$ 丢弃求值器。因此确实有：
\[U^T \circ F^T = T\]
如预期，伴随关系的单位就是单子 $T$ 的单位。

你或许记得伴随关系的余单位通过下述公式产生单子乘法：
\[\mu = R \circ \varepsilon \circ L\]
这是三个自然变换的水平合成，其中两个是分别将 $L$ 映射到 $L$ 和 $R$ 映射到 $R$ 的恒等自然变换。中间的余单位是一个在代数 $(a, f)$ 处分量为 $f$ 的自然变换。

计算分量 $\mu_a$：先将 $\varepsilon$ 与 $F^T$ 水平合成，得到 $F^T a$ 处的分量。由于 $F^T$ 将 $a$ 映射为代数 $(T a, \mu_a)$，而 $\varepsilon$ 选取求值器，最终得到 $\mu_a$。左侧与 $U^T$ 的水平合成不改变结果，因为 $U^T$ 在态射上的作用是平凡的。因此，从伴随关系得到的 $\mu$ 与原单子 $T$ 的 $\mu$ 完全一致。

\section{Kleisli范畴}

我们之前已经见过Kleisli范畴。它是从另一个范畴 $\cat{C}$ 和一个单子(monad) $T$ 构造而来的范畴。我们将这个范畴记作 $\cat{C}_T$。Kleisli范畴 $\cat{C}_T$ 中的对象就是 $\cat{C}$ 的对象，但态射(morphism)有所不同。在Kleisli范畴中从 $a$ 到 $b$ 的态射 $f_{\cat{K}}$ 对应于原范畴中从 $a$ 到 $T b$ 的态射 $f$。我们将这种态射称为从 $a$ 到 $b$ 的Kleisli箭头(Kleisli arrow)。

Kleisli范畴中态射的复合是通过Kleisli箭头的单子复合来定义的。例如，令 $g_{\cat{K}}$ 接续 $f_{\cat{K}}$ 进行复合。在Kleisli范畴中有：
\begin{gather*}
  f_{\cat{K}} \Colon a \to b \\
  g_{\cat{K}} \Colon b \to c
\end{gather*}
这在范畴 $\cat{C}$ 中对应于：
\begin{gather*}
  f \Colon a \to T b \\
  g \Colon b \to T c
\end{gather*}
我们定义复合：
$$h_{\cat{K}} = g_{\cat{K}} \circ f_{\cat{K}}$$
作为 $\cat{C}$ 中的Kleisli箭头：
\begin{align*}
  h & \Colon a \to T c          \\
  h & = \mu \circ (T g) \circ f
\end{align*}
在Haskell中我们会这样编写：

\src{snippet05}

存在一个从 $\cat{C}$ 到 $\cat{C}_T$ 的函子 $F$，它在对象上平凡地作用。对于态射，它将 $\cat{C}$ 中的态射 $f$ 映射为 $\cat{C}_T$ 中通过修饰 $f$ 的返回值生成的Kleisli箭头。给定态射：
$$f \Colon a \to b$$
它会在 $\cat{C}_T$ 中创建对应的Kleisli箭头：
$$\eta \circ f$$
在Haskell中我们会这样编写：

\src{snippet06}

我们也可以定义一个从 $\cat{C}_T$ 回到 $\cat{C}$ 的函子 $G$。它将Kleisli范畴中的对象 $a$ 映射为 $\cat{C}$ 中的对象 $T a$。对于对应Kleisli箭头：
$$f \Colon a \to T b$$
的态射 $f_{\cat{K}}$，其作用是在 $\cat{C}$ 中给出态射：
$$T a \to T b$$
具体方法是对 $f$ 进行提升后再应用 $\mu$：
$$\mu_{T b} \circ T f$$
用Haskell表示法写作：

\begin{snipv}
G f\textsubscript{T} = join . fmap f
\end{snipv}
您可能已认出这是使用 \code{join} 定义单子绑定(bind)的标准形式。

容易看出这两个函子构成了一个伴随(adjunction)：
$$F \dashv G$$
它们的合成 $G \circ F$ 则重现了原始单子 $T$。

这是我们见到的第二个产生相同单子的伴随对。事实上，所有导致 $\cat{C}$ 上相同单子 $T$ 的伴随构成的范畴 $\cat{Adj}(\cat{C}, T)$ 中，Kleisli伴随是该范畴中的始对象(initial object)，而Eilenberg-Moore伴随则是终对象(terminal object)。

\section{余单子的余代数}

对于任意\hyperref[comonads]{余单子}$W$都可以进行类似的构造。我们可以定义与余单子相容的余代数范畴。这些余代数满足以下交换图：

\begin{figure}[H]
  \centering
  \begin{subfigure}
    \centering
    \begin{tikzcd}[column sep=large, row sep=large]
      a \arrow[rd, equal]
      & Wa \arrow[l, "\varepsilon_a"'] \\
      & a \arrow[u, "\mathit{coa}"']
    \end{tikzcd}
  \end{subfigure}%
  \hspace{1cm}
  \begin{subfigure}
    \centering
    \begin{tikzcd}[column sep=large, row sep=large]
      W(Wa)
      & Wa \arrow[l, "W\ \mathit{coa}"'] \\
      Wa \arrow[u, "\delta_a"]
      & a \arrow[u, "\mathit{coa}"] \arrow[l, "\mathit{coa}"']
    \end{tikzcd}
  \end{subfigure}
\end{figure}

\noindent
其中$\mathit{coa}$是载体为$a$的余代数的余评估态射：
$$\mathit{coa} \Colon a \to W a$$
而$\varepsilon$和$\delta$是定义余单子的两个自然变换（在Haskell中其分量被称为\code{extract}和\code{duplicate}）。

从这类余代数范畴到$\cat{C}$存在一个明显的遗忘函子$U^W$，它遗忘余评估结构。我们将考虑其右伴随函子$F^W$。
$$U^W \dashv F^W$$
遗忘函子的右伴随称为余自由函子。$F^W$生成余自由余代数，它将$\cat{C}$中的对象$a$映射到余代数$(W a, \delta_a)$。该伴随对复合产生的函子$U^W \circ F^W$恰好是原始的余单子。

类似地，我们可以构造余Kleisli范畴（co-Kleisli category），通过余Kleisli箭头（co-Kleisli arrows）从对应的伴随关系重新生成余单子结构。

\section{透镜}

让我们回到对透镜的讨论。透镜可以表示为如下余代数：
$$\mathit{coalg}_s \Colon a \to \mathit{Store}\ s\ a$$
对应函子 $\mathit{Store}\ s$:

\src{snippet07}
该余代数也可表示为一对函数：
\begin{align*}
  \mathit{set} & \Colon a \to s \to a \\
  \mathit{get} & \Colon a \to s
\end{align*}
（可将 $a$ 理解为"整体"，$s$ 理解为其"局部"）。用这对函数表示为：
$$\mathit{coalg}_s\ a = \mathit{Store}\ (\mathit{set}\ a)\ (\mathit{get}\ a)$$
其中 $a$ 是类型为 $a$ 的值。注意部分应用的 \code{set} 是函数 $s \to a$。

我们还知道 $\mathit{Store}\ s$ 是一个余单子：

\src{snippet08}
问题是：在什么条件下透镜是该余单子的余代数？第一个一致性条件：
$$\varepsilon_a \circ \mathit{coalg} = \idarrow[a]$$
转化为：
$$\mathit{set}\ a\ (\mathit{get}\ a) = a$$
这就是透镜定律之一，表明若将结构 $a$ 的字段设为其当前值，则结构保持不变。

第二个条件：
$$\mathit{fmap}\ \mathit{coalg} \circ \mathit{coalg} = \delta_a \circ \mathit{coalg}$$
需要更多推导。首先回顾 \code{Store} 函子的 \code{fmap} 定义：

\src{snippet09}
对 \code{coalg} 的结果应用 \code{fmap coalg} 得到：

\src{snippet10}
另一方面，对 \code{coalg} 的结果应用 \code{duplicate} 得到：

\src{snippet11}
要使这两个表达式相等，当作用于任意 \code{s} 时，\code{Store} 下的两个函数必须相等：

\src{snippet12}
展开 \code{coalg} 后得到：

\src{snippet13}
这对应剩余的两个透镜定律。第一个：

\src{snippet14}
表明连续两次设置字段值等同于仅设置一次。第二个定律：

\src{snippet15}
表明获取被设置为 $s$ 的字段值会返回 $s$。

换句话说，行为良好的透镜确实是 \code{Store} 函子的余单子余代数。

\section{挑战}

\begin{enumerate}
  \tightlist
  \item
        自由函子 $F \Colon C \to C^T$ 在态射上的作用是什么？提示：利用单子乘法(monadic $\mu$)的自然性条件(naturality condition)。
  \item
        定义伴随对(adjunction)：
        \[U^W \dashv F^W\]
  \item
        证明上述伴随对重现原始的余单子(comonad)。
\end{enumerate}